\scanpage{58}
\noindent\hfill atkin page 2\mspage{2}\par
\begingroup
\csname tagsleft@true\endcsname

2. \handwritten{\uline}{The} \handwritten{\uline}{full} \handwritten{\uline}{modular} \handwritten{\uline}{group}. We consider the modular group $\handwritten{\underline}{\handwritten{\Gamma}}$ consisting of all linear fractional transformations
\[
z\ \rightarrow\ Vz=(az+b)/(cz+d)
\quad \text{(with } \handwritten{\underline}{a,b,c,d},\text{rational integers and }\handwritten{\underline}{ad-bc\mathrel{\handwritten{=}}}1)
\]
acting on the upper half $\handwritten{\underline}{z}$-plane $\handwritten{\underline}{H}$,and write
\[
\handwritten{\underline}{x}=\handwritten{\underline}{e}^{2\handwritten{\underline}{\handwritten{\pi}}\handwritten{\underline}{iz}} \qquad \unsure{.}
(\,|x|<1\tikz[remember picture,baseline]{\node[inner sep=0pt](punctfiftyeight){$\deleted{\cdot}\mathpunct{,}$};}\quad \mathrm{Im}\ \handwritten{\underline}{z}>0\,)\ .
\]
\reversemarginpar\margin{\handwritten{\unsure{.} pi}}\normalmarginpar
\margin{\handwritten{\tikz[remember picture,baseline]{\node[inner sep=0pt](delefiftyeight){\tikz[baseline=-.4ex,x=.85em,y=.85em]{\draw (.03,-.7) .. controls (.36,-.06) and (.83,.64) .. (.59,.63) .. controls (.18,.63) and (.16,.06) .. (.45,.22) .. controls (.61,.35) and (.32,.56) .. (.46,.68) -- (.93,1.12);}};}\tikz[remember picture,overlay]{\draw[->] (punctfiftyeight.north east) to[bend left=6] (delefiftyeight.west);}}}

We define a \handwritten{\uline}{modular} \handwritten{\uline}{form} $\handwritten{\underline}{F(z)}$ of dimension $\handwritten{\underline}{k}$(which we take to be an even integer throughout this paper) by
\[
\handwritten{\underline}{F(Vz)=(cz+d)^{-k}F(z)\quad \text{for all } V \handwritten{\in}\handwritten{\underline}{\handwritten{\Gamma}}}\quad ,
\tag*{(M1)}
\]
\[
\parbox[t]{.82\linewidth}{\raggedright \handwritten{\uline}{$F(z)$ is meromorphic in the fundamental domain $D=I+B$ of $\handwritten{\underline}{\handwritten{\Gamma}}$} ,}
\tag*{(M2)}
\]
\handwritten{\uline}{where $I$ is that part of $D$ inside $H$,and $B$(the cusps of $D$) is that part of $D$ on the boundary of $H$; here we may take $B$ as the single point $z=i\handwritten{\infty}$.}

The conditions (M1) and (M2) seem insufficient to produce much number theory,and we therefore define further restrictions as follows:
\[
\text{an \handwritten{\uline}{entire} form is regular in } I\quad ,
\tag*{(M3a)}
\]
\[
\text{a \handwritten{\uline}{regular} form is regular in } D=I+B\quad ,
\tag*{(M3b)}
\]
\[
\text{a \handwritten{\uline}{cuspform} is regular in } I \text{ and zero in } B.
\tag*{(M3c)}
\]
Clearly a cuspform is a regular form,and a regular form is an entire form. We denote the vector spaces of entire,regular,and cusp-,forms(for the given dimension $\handwritten{\underline}{k}$) by $\handwritten{\underline}{C},\handwritten{\underline}{C}^{\handwritten{+}}$,and $\handwritten{\underline}{C}^{\handwritten{0}}\handwritten{,}$respectively. The last two \tikz[baseline=(old.base)]{\node[inner sep=0pt](old){\deleted{\unsure{aure|aare}}};\draw (old.center) ellipse [x radius=1em,y radius=.65em];}\inserted{\handwritten{\tikz[baseline=(new.base)]{\node[inner sep=0pt](new){are};\draw (new.center) ellipse [x radius=.9em,y radius=.65em];}}} finite-dimensional($\handwritten{\underline}{C}$ is not). Since $\handwritten{\underline}{C}^{\handwritten{+}}$ is the direct sum of $\handwritten{\underline}{C}^{\handwritten{0}}$ and a well-defined Eisenstein series,we ignore $\handwritten{\underline}{C}^{\handwritten{+}}$ from now on.

By the \handwritten{\uline}{Fourier} \handwritten{\uline}{series} (FS)\unsure{-} of a modular form we understand its expansion in powers of $\handwritten{\underline}{x}$,the local variable at the cu\deleted{p}sp $\handwritten{\underline}{z}=\handwritten{\underline}{i}\handwritten{\infty}$ $\handwritten{\vdots}$
\[
F(z)=F(x)=\handwritten{\sum_{n=N}^{\infty}}a(n)\,x^{n}\qquad (N>-\handwritten{\infty})\qquad ,
\]
and we call \tikz[remember picture,baseline=(oldcoeff.base)]{\node[inner sep=0pt](oldcoeff){\deleted{\illegible{1}}};}\reversemarginpar\margin{\handwritten{\tikz[remember picture,baseline]{\node[inner sep=0pt](deletwo){\tikz[baseline=-.4ex,x=.85em,y=.85em]{\draw (.03,-.7) .. controls (.36,-.06) and (.83,.64) .. (.59,.63) .. controls (.18,.63) and (.16,.06) .. (.45,.22) .. controls (.61,.35) and (.32,.56) .. (.46,.68) -- (.93,1.12);}};}\tikz[remember picture,overlay]{\draw (oldcoeff.north) to[bend right=7] (deletwo.east);}}}\normalmarginpar $\handwritten{\underline}{a(n)}$ the Fourier coefficients. Roughly speaking, the number theory associated with modular theory arises from the Fourier coefficients of entire modular forms. Indeed,prior to Ramanujan and his work on congruences and(with Hardy)asymptotic formulae for $\handwritten{\underline}{p(n)}$,one could read ``regular''for ``entire'' in the last sentence. The author does not know of any number theory associated with modular forms which have poles inside $\handwritten{\underline}{H}$.
\endgroup

% archive sheet number: 57
